\documentclass[11pt]{article}
\usepackage{amsmath}
\usepackage{amssymb}
\usepackage{geometry}
\geometry{margin=1in}

\title{STAT 17 Discussion Notes \\ Week 9: Hypothesis Testing}
\author{Xiao Zhang}
\date{March 2026}

\begin{document}

\maketitle

\section{Goal of Hypothesis Testing}

Hypothesis testing is a statistical method used to determine whether there is enough evidence in sample data to support a claim about a population parameter.

Examples of questions answered using hypothesis tests:

\begin{itemize}
\item Does a new drug have fewer side effects?
\item Has public opinion changed over time?
\item Is the average body temperature different from 98.6°F?
\end{itemize}

A hypothesis test uses sample data to decide whether a population parameter is likely equal to a hypothesized value.

\section{The Five Steps of a Hypothesis Test}

\begin{enumerate}

\item \textbf{State the Hypotheses}

Two hypotheses are always defined:

\begin{itemize}
\item Null hypothesis ($H_0$): the status quo or assumed value
\item Alternative hypothesis ($H_1$ or $H_a$): the claim we want evidence for
\end{itemize}

Three possible forms:

\[
H_0: \text{parameter} = \text{value}
\]

Left-tailed:
\[
H_1: \text{parameter} < \text{value}
\]

Right-tailed:
\[
H_1: \text{parameter} > \text{value}
\]

Two-tailed:
\[
H_1: \text{parameter} \neq \text{value}
\]

\item \textbf{Check Conditions and Choose Significance Level}

Common significance levels:

\[
\alpha = 0.01, \quad 0.05, \quad 0.10
\]

$\alpha$ represents the probability of making a Type I error.

\item \textbf{Calculate the Test Statistic}

The test statistic measures how far the sample statistic is from the hypothesized value in standard errors.

\item \textbf{Find the p-value}

The p-value is the probability of observing the sample result (or something more extreme) assuming $H_0$ is true.

\item \textbf{Draw a Conclusion}

\begin{itemize}
\item If $p$-value $\leq \alpha$: Reject $H_0$
\item If $p$-value $> \alpha$: Do Not Reject $H_0$
\end{itemize}

Conclusions should always be written in terms of the \textbf{alternative hypothesis}.

\end{enumerate}

\section{Hypothesis Test for a Population Proportion}

Used when the parameter is a population proportion $p$.

\subsection{Conditions}

\begin{itemize}
\item Random sample or randomized experiment
\item Sample size less than 5\% of population
\item Large sample condition:
\[
np_0(1-p_0) \ge 10
\]
\end{itemize}

\subsection{Test Statistic}

\[
z =
\frac{\hat{p}-p_0}
{\sqrt{\frac{p_0(1-p_0)}{n}}}
\]

\subsection{p-value}

Left-tailed test:
\[
p\text{-value}=P(Z<z)
\]

Right-tailed test:
\[
p\text{-value}=P(Z>z)
\]

Two-tailed test:
\[
p\text{-value}=2P(Z>|z|)
\]

\section{Hypothesis Test for a Population Mean (Case 1: $\sigma^2$ is unknown)}

Used when the parameter is a population mean $\mu$, $\sigma$ .

\subsection{Conditions}

\begin{itemize}
\item Random sample
\item Sample size less than 5\% of population
\item Data from a normal population OR $n \ge 30$
\end{itemize}

\subsection{Test Statistic}

\[
t =
\frac{\bar{x}-\mu_0}
{s/\sqrt{n}}
\]

\begin{itemize}
\item Follows a $t$ distribution
\item Degrees of freedom: $df = n-1$
\end{itemize}

\subsection{p-value}

Left-tailed:
\[
p\text{-value}=P(t<t_0)
\]

Right-tailed:
\[
p\text{-value}=P(t>t_0)
\]

Two-tailed:
\[
p\text{-value}=2P(t>|t_0|)
\]


\section{Hypothesis Test for a Population Mean (Case 2)}

Used when the parameter is a population mean $\mu$.

\subsection{Conditions}

\begin{itemize}
\item Random sample
\item Sample size less than 5\% of population
\item Data from a normal population OR $n \ge 30$
\end{itemize}

\subsection{Test Statistic}

\[
z =
\frac{\bar{x}-\mu_0}
{\sigma/\sqrt{n}}
\]

\subsection{p-value}

Left-tailed:
\[
p\text{-value}=P(Z<z_0)
\]

Right-tailed:
\[
p\text{-value}=P(Z>z_0)
\]

Two-tailed:
\[
p\text{-value}=2P(Z>|z_0|)
\]


\section{Relationship Between Confidence Intervals and Hypothesis Tests}

For a two-sided test:

\begin{itemize}
\item If the hypothesized value is \textbf{outside} the confidence interval $\Rightarrow$ Reject $H_0$
\item If the hypothesized value is \textbf{inside} the confidence interval $\Rightarrow$ Do Not Reject $H_0$
\end{itemize}

Example:

A 95\% confidence interval that contains the null value means the hypothesis test at $\alpha = 0.05$ will not reject $H_0$.

\section{Type I and Type II Errors}

Because hypothesis tests use sample data, mistakes are possible.

\subsection{Type I Error}

Rejecting $H_0$ when $H_0$ is actually true.

\[
P(\text{Type I error}) = \alpha
\]

Example: Concluding a drug has an effect when it actually does not.

\subsection{Type II Error}

Failing to reject $H_0$ when $H_1$ is actually true.

\[
P(\text{Type II error}) = \beta
\]

Example: Concluding a drug has no effect when it actually works.

\subsection{Important Relationship}

Reducing $\alpha$ increases $\beta$.

\section{Key Summary}

\textbf{Hypothesis testing process}

\begin{enumerate}
\item State $H_0$ and $H_1$
\item Check conditions
\item Compute test statistic
\item Find p-value
\item Compare with $\alpha$ and conclude
\end{enumerate}

\end{document}

